\newcommand{\touche}[1]{{
\begin{picture}(12,12)(0, 2)
\put(0,0){\framebox(12,12)[lb]{}}
\put(1.5,2){\makebox(0,0)[lb]{\tiny #1}}
\end{picture}\hspace{0.3ex}}}


\section{Le robot} 
%\label{sec:2:le-robot} 

Pour réaliser nos dessins géométriques sur ordinateur nous allons piloter
un petit robot. Son nom est \robot\@\footnote{\ La calligraphie de son nom
ne laisse aucun doute sur son caractère humanoïde et sympathique. 
}. 
Notre \robot\ porte un crayon qu'il peut le\-ver ou 
bais\-ser pen\-dant ses dé\-pla\-ce\-ments sur l'écran de l'ordinateur.  
Quand le \robot\ se 
déplace avec son crayon baissé, il laisse une trace sur l'écran 
comme le ferait un escargot sur une feuille. Quand il se déplace avec son 
crayon levé, il ne laisse aucune trace de son passage.  Pour dessiner,
le \robot\ doit suivre exactement les lignes du dessin avec son crayon 
baissé.
Il peut également aller d'un point à un autre du dessin qui ne sont pas 
reliés par un trait en levant son crayon. Pour guider le
\robot\ nous utilisons le clavier de l'ordinateur qui sert de télécommande.
Nous pouvons donc piloter le \robot\  comme s'il s'agissait d'une petite voiture 
conduite à distance par ra\-dio\-com\-man\-de.

\smallskip

Sur l'écran de l'ordinateur notre \robot\ est représenté par une 
petite flèche \robotCbN\ qui indique dans quelle direction de la
rose des vents il est tourné (voir la figure 
\ref{fig:rose-des-vents-8d}, page \pageref{fig:rose-des-vents-8d}). Il peut s'orienter
au nord \robotCbN\/, au nord-est \robotCbNE, à l'est \robotCbE\/, au
sud-est \robotCbSE\/, au sud \robotCbS\/, au sud-ouest \robotCbSW\/, à
l'ouest \robotCbW\ ou au nord-ouest \robotCbNW\/.
La flè\-che est pleine quand le \robot\ a baissé son crayon et
évidée
(\robotClN\/, \robotClNE\/, \robotClE\/, \robotClSE\/, \robotClS\/,
\robotClSW\/,  \robotClW\ ou \robotClNW\/) quand le crayon est levé.
\begin{figure}[htbp]
 \includehvspacePostScript[6cm,5.6cm]{rose-des-vents-8d}{0.95}{0.95}
 \caption{Rose des vents à huit directions}\label{fig:rose-des-vents-8d}
\end{figure}

\smallskip

Comme les trains sur leurs rails, notre \robot\ ne peut pas aller 
n'importe où sur l'écran. Il doit impérativement se déplacer sur 
les traits ou arcs joignant deux points successifs d'une grille dont les dimensions sont 
fixées par la taille de l'écran de l'ordinateur. Une petite partie de cette grille
est  représentée \hbox{ci-dessous} à la figure \ref{fig:grille}.
\index{Grille du \protect\robot}%
\begin{figure}[htbp]
 \includehvspacePostScript[\textwidth,4.0cm]{grille}{0.95}{0.95}
 \caption{Grille et points d'arrêt du \protect\robot}
 \label{fig:grille}
\end{figure}

\noindent Cette grille imaginaire correspond au papier quadrillé que nous avons utilisé pour
faire nos dessins à la main, sauf que les carreaux ne sont pas visibles sur 
l'écran de l'ordinateur. 
Après un déplacement, le \robot\ s'arrête toujours à l'intersection
d'une ligne horizontale et d'une ligne verticale de la grille imaginaire. La figure
 \ref{fig:traject} (page \pageref{fig:traject}) montre toutes les trajectoires possibles
du \robot\ entre deux points d'arrêt quelconques sur la grille.
\begin{figure}[htbp]
 \includevspacePostScript[5.4cm]{dessincomplexe}{0.95}{0.95}
 \caption{Trajectoires du \protect\robot\ (grossissement $\times$\ $4,5$)}
 \label{fig:traject}
\end{figure}
\begin{figure}[htbp]
 \includehvspacePostScript[42mm,2.1cm]{robot-arrete}{0.95}{0.95}
 \caption{Le \protect\robot\ arrêté (grossissement $\times$\ $2,5$)}
 \label{fig:robot-arrete}
\end{figure}
La figure \ref{fig:robot-arrete} \hbox{(page \pageref{fig:robot-arrete})} 
décrit la position
exacte du \robot\ quand il s'arrête sur la grille. Il place la pointe de
son crayon située au milieu et à l'arrière de la flèche exactement sur
un point d'arrêt de la grille.


\section{Pilotage du robot} 

Pour piloter le \robot\ sur l'ordinateur, il faut ``{\em exécuter le programme}
{\idt Des\-si\-ner}'' qui se trouve sur la disquette
d'accompagnement de ce livre. La façon d'exécuter ce programme 
{\tt Des\-si\-ner} dépend de l'ordinateur utilisé%
\footnote{\ Introduire dans un lecteur la disquette 
d'accompagnement de ce livre puis sur compatible IBM PC taper 
{\tt Des\-si\-ner} ou sur Macintosh cliquer deux fois
rapidement sur l'icône de l'application {\tt Des\-si\-ner}.}. 

\subsection{Position initiale du robot}

Quand l'exécution du programme {\idt Des\-si\-ner}
commence, l'écran 
de l'ordinateur se présente comme on le voit en réduction à la figure 
\ref{fig:ecran-initial-1}  (page \pageref{fig:ecran-initial-1}).
\begin{figure}[htbp]
 \includevspacePostScript[7.6cm]{ecran-initial-1}{0.95}{0.95}
 \caption{Écran de l'ordinateur au début}
 \centering{de l'exécution du programme {\idt Des\-si\-ner}}
 \label{fig:ecran-initial-1}
\end{figure}
Au départ le \robot\ se trouve donc arrêté au milieu de la grille
imaginaire, crayon baissé, or\-ienté vers le nord \robotCbN\/.
Pour piloter le \robot, nous utilisons les touches du clavier de 
l'ordinateur qui sert de télécommande. Nous allons maintenant
voir sur quelles touches appuyer pour lever ou baisser le crayon, 
avancer, pivoter sur place, tourner, etc.

\subsection{Manœuvrer le crayon du robot}

 En tapant sur la touche marquée \touche{L} du
clavier, le \robot\ lève son crayon \robotClN\/. En tapant sur la touche 
\touche{B}, il le baisse \robotCbN\/. Faute de disposer de beaucoup de place sur l'écran 
de l'ordinateur pour
permettre de rappeler en détail le mode d'emploi du \robot, nous utilisons
des abréviations mnémotechniques% 
%
\footnote{\ Les abréviations mnémotechniques permettent d'aider la mémoire par association d'idées.}
placées en bas de l'écran comme on peut
le voir à la figure 
\ref{fig:ecran-initial-1} ci-dessus. 
Par \hbox{exemple~:}
\begin{list}{$\scriptscriptstyle\bullet$}{
}
\item {\bf L}$:${\idt lc} signifie ``la touche \touche{L} permet de {\underline l}ever le \hbox{{\underline c}rayon''$\:$;}
\item {\bf B}$:${\idt bc} signifie ``la touche \touche{B} permet de {\underline b}aisser le \hbox{{\underline c}rayon''.}
\end{list}

L'effet des touches \touche{L} et \touche{B} de manœuvre du crayon du 
\robot\ est résumé dans la table \ref{tab:crayon} (page \pageref{tab:crayon}).

\begin{table}[htbp]
\begin{center}
\begin{tabular}{@{}|c|c|c|c|c|c|c|c|c|c|@{}} \hline
\touche{L} ({\idt lc}) & Avant   & \robotCbN  & \robotCbNE & \robotCbE & \robotCbSE & \robotCbS  & \robotCbSW
& \robotCbW & \robotCbNW \\ \cline{2-10}
                     & Après & \robotClN  & \robotClNE & \robotClE & \robotClSE  & \robotClS  
& \robotClSW & \robotClW  & \robotClNW  \\ \hline\hline
\touche{B} ({\idt bc}) & Avant  & \robotClN  & \robotClNE & \robotClE & \robotClSE & \robotClS  & \robotClSW  
& \robotClW  & \robotClNW  \\ \cline{2-10}
                         & Après & \robotCbN & \robotCbNE & \robotCbE  & \robotCbSE   & \robotCbS 
& \robotCbSW & \robotCbW  & \robotCbNW \\ \hline
\end{tabular}
\end{center}
\caption{Manœuvres du crayon du \protect\robot}
\label{tab:crayon}
\end{table}



\subsection{Faire pivoter le robot sur place d'un quart de tour}


La touche \touche{$-$} fait \underline{p}ivoter le \robot\ sur place d'un 
\underline{q}uart de \underline{t}our à \underline{g}auche ({\idt pqtg} en abrégé).
Par exemple si le \robot\ est orienté dans la direction nord \robotCbN\/, la touche \touche{$-$} le fait
pivoter sur place dans la direction ouest \robotCbW\/. S'il est orienté à l'est \robotCbE\/, 
il se place dans la direction nord \robotCbN\/. 
La touche \touche{$+$} fait \underline{p}ivoter le \robot\ sur place d'un 
\underline{q}uart de \underline{t}our à \underline{d}roite ({\idt pqtd} en abrégé).
 
Quand on parle du côté gauche du \robot, il faut s'imaginer à
la place du \robot, le regard tourné dans la direction de 
la flèche. Par exemple, la gauche de \robotCbE\ est le haut de cette page tandis 
que la gauche de \robotCbS\ est à droite de cette \hbox{feuille$\:$!}
 
L'effet des commandes \touche{$-$} ({\idt pqtg})  et \touche{$+$} ({\idt pqtd}) 
est résumé dans la table ci-dessous dans le cas où le crayon est baissé. L'effet est
le même quand le crayon est levé.
\begin{table}[htbp]
\begin{center}
\begin{tabular}{@{}|c|c|c|c|c|c|c|c|c|c|@{}} \hline
\touche{$-$} ({\idt pqtg}) & Avant   & \robotCbN & \robotCbNE & \robotCbE & \robotCbSE & \robotCbS & \robotCbSW & \robotCbW & \robotCbNW \\ \cline{2-10}
                         & Après & \robotCbW & \robotCbNW & \robotCbN & \robotCbNE & \robotCbE & \robotCbSE & \robotCbS & \robotCbSW \\ \hline \hline
\touche{$+$} ({\idt pqtd}) & Avant   & \robotCbN & \robotCbNE & \robotCbE & \robotCbSE & \robotCbS & \robotCbSW & \robotCbW & \robotCbNW \\ \cline{2-10}
                         & Après & \robotCbE & \robotCbSE & \robotCbS & \robotCbSW & \robotCbW & \robotCbNW & \robotCbN & \robotCbNE \\ \hline
\end{tabular}
\end{center}
\caption{Pivotements du \protect\robot\ d'un quart de tour sur lui-même}
%\label{tab:pqtdg}
\end{table}

\subsection{Faire pivoter le robot sur place d'un huitième de tour}

La touche \touche{G} fait \underline{p}ivoter le \robot\ sur place d'un 
\underline{h}uitième de \underline{t}our à \underline{g}auche ({\idt phtg} en abrégé).
Par exemple, s'il est dans la direction nord \robotCbN\/, le \robot\ se place dans la direction
nord-ouest \robotCbNW\/. En continuant à taper sur la touche \touche{G}, le \robot\
tourne à chaque fois d'un huitième de tour en passant par les positions successives
\robotCbW\ puis \robotCbSW\/, \robotCbS\/, \robotCbSE\/, \robotCbE\/, \robotCbNE\/, \ldots
La touche \touche{D} fait \underline{p}ivoter le \robot\ sur place d'un 
\underline{h}uitième de \underline{t}our à \underline{d}roite ({\idt phtd} en abrégé).
L'effet des touches  \touche{G} ({\idt phtg})  et \touche{D} ({\idt phtd}) est résumé 
\hbox{ci-dessous~:}
\begin{table}[htbp]
\begin{center}
\begin{tabular}{@{}|c|c|c|c|c|c|c|c|c|c|@{}} \hline
\touche{G} ({\idt phtg}) & Avant    & \robotCbN  & \robotCbNW & \robotCbW  & \robotCbSW & \robotCbS  & \robotCbSE & \robotCbE  & \robotCbNE \\ \cline{2-10}
                       & Après  & \robotCbNW & \robotCbW  & \robotCbSW & \robotCbS  & \robotCbSE &\robotCbE   & \robotCbNE & \robotCbN  \\ \hline \hline
\touche{D} ({\idt phtd}) & Avant    & \robotCbN  & \robotCbNE & \robotCbE  & \robotCbSE & \robotCbS  & \robotCbSW & \robotCbW  & \robotCbNW \\ \cline{2-10}
                       & Après & \robotCbNE & \robotCbE  & \robotCbSE & \robotCbS  & \robotCbSW & \robotCbW  & \robotCbNW & \robotCbN  \\ \hline
\end{tabular}
\end{center}
\caption{Pivotements du \protect\robot\ d'un huitième de tour sur lui-même}
%\label{tab:phtgd}
\end{table}

\subsection{Faire avancer le robot}

La touche \touche{A} fait \underline{av}ancer le \robot\ tout droit, de la taille d'un 
côté d'un carreau du quadrillage. 
L'abréviation mnémotechnique de cette commande du \robot\ est {\idt av}.
Si le crayon est baissé, ce déplacement permet de tracer un petit trait. Si le crayon 
est levé, ce déplacement ne laisse aucune trace sur le dessin. Dans les deux cas le
\robot\ avance tout droit dans la direction où il se trouve sans 
changer \hbox{d'orientation~:}
\begin{table}[htbp]
\begin{center}
\input ROBOT-AVPICT2.46x2.46.tex
\input ROBOT+AVPICT2.46x2.46.tex
\input ROBOT-VGPICT2.46x2.46.tex 
\input ROBOT+VGPICT2.46x2.46.tex 
\input ROBOT-VDPICT2.46x2.46.tex
\input ROBOT+VDPICT2.46x2.46.tex
\begin{tabular}{@{}|c|c|c|c|c|c|c|c|c|c|@{}} \hline
\touche{A} ({\idt av}) & Avant   & \robotCbNmav & \robotCbNEmav & \robotCbEmav & \robotCbSEmav & \robotCbSmav & \robotCbSWmav & \robotCbWmav & \robotCbNWmav \\ \cline{2-10}
                     & Après & \robotCbNpav & \robotCbNEpav & \robotCbEpav & \robotCbSEpav & \robotCbSpav & \robotCbSWpav & \robotCbWpav & \robotCbNWpav \\ \hline 
\end{tabular}
\end{center}
\caption{Avancées du \protect\robot\ dans les huit directions}
%\label{tab:av-8d}
\end{table}

\begin{exemple}[Frise crantée]{\rm
 En tapant successivement sur les touches \touche{A} ({\idt av}), \touche{$+$} ({\idt pqtd}),  \touche{A} ({\idt av}),
\touche{$+$} ({\idt pqtd}), \touche{A} ({\idt av}), \touche{$-$} ({\idt pqtg}), \touche{A} ({\idt av}), 
\touche{$-$} ({\idt pqtg}), \touche{A} ({\idt av}), \touche{$+$} ({\idt pqtd}),  \touche{A} ({\idt av}),
\touche{$+$} ({\idt pqtd}), \touche{A} ({\idt av}), \touche{$-$} ({\idt pqtg}), \touche{A} ({\idt av}), 
\touche{$-$} ({\idt pqtg}), \ldots, on ob\-tient la frise \hbox{suivante~:}

\includevspacePostScript[0.7cm]{frise-crantee}{0.95}{0.95}
\noindent $\Box$ }\end{exemple}

\begin{exemple}[Frise crénelée]{\rm
 En tapant successivement sur les touches
\touche{D} ({\idt phtd}), \touche{D} ({\idt phtd}), \touche{A} ({\idt av}), \touche{D} ({\idt phtd}), 
\touche{A} ({\idt av}), \touche{G} ({\idt phtg}), \touche{A} ({\idt av}), \touche{G} ({\idt phtg}), 
\touche{A} ({\idt av}), \touche{D} ({\idt phtd}), \touche{A} ({\idt av}), \touche{D} ({\idt phtd}), 
\touche{A} ({\idt av}), \ldots, on ob\-tient la frise \hbox{suivante~:}

\includevspacePostScript[0.7cm]{frise-crenelee}{0.95}{0.95}
\noindent $\Box$ }\end{exemple}

\subsection{Faire virer le robot}

La touche \touche{(} \vspace{-1.0mm} permet de faire {\underline v}irer le \robot\ à {\underline g}auche 
({\idt vg} en abrégé) tandis que la
touche \touche{)} permet de le faire {\underline v}irer 
à {\underline d}roite ({\idt vd}).
Ces touches ont été
choisies parce que les parenthèses évoquent l'idée de virage et que
la parenthèse gauche \touche{(} qui se trouve à gauche sur le clavier 
correspond au virage à gauche tandis que la
parenthèse droite \touche{)} qui correspond au virage à droite se trouve à droite sur le 
clavier. Si le crayon est baissé, ce déplacement permet de tracer
un petit arc de cercle. Si le crayon est levé, ce déplacement ne laisse 
aucune trace sur le dessin. Dans les deux cas le \robot\ se déplace sur la
grille et change d'orientation comme indiqué dans la table 
\hbox{ci-dessous~:}
\begin{table}[htbp]
\begin{center}
\input ROBOT-AVPICT2.46x2.46.tex
\input ROBOT+AVPICT2.46x2.46.tex
\input ROBOT-VGPICT2.46x2.46.tex 
\input ROBOT+VGPICT2.46x2.46.tex 
\input ROBOT-VDPICT2.46x2.46.tex
\input ROBOT+VDPICT2.46x2.46.tex
\begin{tabular}{@{}|c|c|c|c|c|c|c|c|c|c|@{}} \hline
\touche{(} ({\idt vg}) & Avant & \robotCbNmvg & \robotCbNEmvg & \robotCbEmvg & \robotCbSEmvg & \robotCbSmvg & \robotCbSWmvg & \robotCbWmvg & \robotCbNWmvg \\\cline{2-10}
                       & Après & \robotCbNpvg & \robotCbNEpvg & \robotCbEpvg & \robotCbSEpvg & \robotCbSpvg & \robotCbSWpvg & \robotCbWpvg & \robotCbNWpvg \\\hline \hline
\touche{)} ({\idt vd}) & Avant & \robotCbNmvd & \robotCbNEmvd & \robotCbEmvd & \robotCbSEmvd & \robotCbSmvd & \robotCbSWmvd & \robotCbWmvd & \robotCbNWmvd \\\cline{2-10}
                       & Après & \robotCbNpvd & \robotCbNEpvd & \robotCbEpvd & \robotCbSEpvd & \robotCbSpvd & \robotCbSWpvd & \robotCbWpvd & \robotCbNWpvd \\\hline
\end{tabular}
\end{center}
 \caption{Virages du \protect\robot}
%\label{tab:()-vgd}
\end{table}

\begin{exemple}[Frise d'accolades]{\rm
 En tapant successivement sur les tou\-ches 
 \touche{)} ({\idt vd}), \touche{(} ({\idt vg}), 
 \touche{$+$} ({\idt pqtd}), \touche{$+$} ({\idt pqtd}),
 \touche{(} ({\idt vg}), \touche{)} ({\idt vd}), 
 \touche{$+$} ({\idt pqtd}), \touche{$+$} ({\idt pqtd}),
 \touche{)} ({\idt vd}), \touche{(} ({\idt vg}),
 \touche{$+$} ({\idt pqtd}), \touche{$+$} ({\idt pqtd}),
\ldots, on ob\-tient la frise \hbox{suivante~:}

\includevspacePostScript[1.0cm]{frise-accolades}{0.95}{0.95} 
\noindent $\Box$}\end{exemple}

\subsection{Impossibilité de sortir du cadre}

En tapant sur la touche  \touche{A} plusieurs fois de suite le \robot\ \robotCbN\ avance tout droit
en traçant un trait jusqu'à rencontrer le bord du cadre qui délimite son champ d'action.  
Si l'on continue à taper sur la touche
\touche{A} pour essayer de le faire sortir du cadre c'est une {\em erreur} 
en ce sens que
le \robotErr\ refuse de le faire et émet un 
\raisebox{0.25ex}{$\scriptscriptstyle \boldmath\ll$}~%
bip-bip%
~\raisebox{0.25ex}{$\scriptscriptstyle \boldmath\gg$}
de 
protestation\@\footnote{\ La calligraphie de son nom
ne laisse aucun doute sur son mécontentement. 
}$\:$! 

\subsection{Corriger les erreurs de pilotage}

\index{Erreur} Il arrive souvent que l'on fasse des erreurs de pilotage du \robotErr. Par
exemple, au lieu de le faire tourner à droite avec une 
parenthèse droite \touche{)} ({\idt vd}), on tape une parenthèse gauche 
\touche{(} ({\idt vg}) et le \robotErr\ tourne
à gauche$\:$! La touche d'{\em effacement} marquée d'une flèche 
orientée à gauche \index{Touche d'effacement}
\hspace{0.3ex}\framebox[11.mm][l]{$\leftarrow$~~~}\hspace{0.3ex}
permet de corriger les erreurs en revenant en arrière et en effaçant
les derniers tracés dans l'ordre inverse où ils ont été
dessinés par le \robot.

\subsection{Terminer et reproduire le dessin}

Pendant l'exécution du programme {\idt Des\-si\-ner}, 
il est possible,  à chaque ins\-tant, de
terminer le dessin en tapant sur la touche de retour à la ligne 
\hbox{\framebox[11.mm][l]{$\hookleftarrow$~~~}\hspace{0.3ex},} également
marquée {\small\tt Entrée} ou en anglais {\small\tt Return} 
ou {\small\tt Ret} sur certains  ordinateurs. Ceci est rappelé 
en bas à droite de l'écran de l'ordinateur par un symbole 
$\hookleftarrow$ qui clignote. 

Le dessin qui vient d'être terminé n'est pas perdu, il est
possible de le reproduire et de l'im\-pri\-mer%
\footnote{\ La façon de procéder, qui dépend de l'ordinateur utilisé,
est expliquée sur la disquette d'accompagnement du livre.}.
 

\subsection{Dessiner un point, une croix, la grille des points d'arrêt du robot}

La touche \touche{P} ({\idt dp}) permet de \underline{d}essiner un 
\underline{p}oint, sans déplacer le \robot. On utilise la
touche \touche{C} ({\idt dc}) pour \underline{d}essiner une petite
\underline{c}roix {\tt +}. La touche \touche{.} ({\idt dg})
permet de \underline{d}essiner la \underline{g}rille des points 
d'arrêt du \robot\ comme on la voit par exemple sur la figure
\ref{fig:transports} (page \pageref{fig:transports}).

\subsection{Dessiner en couleur}

Quand on dispose d'un ordinateur avec écran couleur, le \robot\ a la
possibilité de dessiner avec des crayons de différentes couleurs.
Pendant l'exé\-cu\-tion du programme {\idt Des\-si\-ner}, 
la couleur du crayon utilisé par le \robot\ figure dans un petit carré 
placé en bas à droite de l'écran (voir la figure \ref{fig:ecran-initial-1}
page \pageref{fig:ecran-initial-1}). Pour changer la couleur du crayon, 
il suffit de taper une ou plusieurs fois sur la touche \touche{*} pour faire
défiler toutes les couleurs possibles les unes
après les autres.

\subsection{Déplacer rapidement le robot avec la souris ou les flèches 
de défilement}

Pour déplacer le \robot\ sans rien dessiner, on peut lever le crayon
\touche{L} ({\idt lc}) puis avancer une ou plusieurs fois 
\touche{A} ({\idt av}) en changeant éventuellement de direction
(\touche{D} ({\idt phtd}) ou \touche{G} ({\idt phtg})) et enfin baisser
le crayon \touche{B} ({\idt bc}). Pour aller plus vite, on peut utiliser
la souris pour déplacer le \robot\ sur un point quelconque de la grille 
des points d'arrêt sans rien tracer. 
Sur les ordinateurs ne disposant pas 
de souris, le même effet peut être obtenu avec les
flèches de défilement 
vers le haut 
\begin{picture}(12,12)(0, 2)
\put(0,0){\framebox(12,12)[l]{$\scriptstyle\uparrow$}}
\end{picture}\hspace{0.3ex},
vers le bas 
\begin{picture}(12,12)(0, 2)
\put(0,0){\framebox(12,12)[l]{$\scriptstyle\downarrow$}}
\end{picture}\hspace{0.3ex},
à gauche 
\begin{picture}(12,12)(0, 2)
\put(0,0){\framebox(12,12)[bl]{$\scriptstyle\leftarrow$}}
\end{picture}\hspace{0.3ex},
ou à droite 
\begin{picture}(12,12)(0, 2)
\put(0,0){\framebox(12,12)[bl]{$\scriptstyle\rightarrow$}}
\end{picture}\hspace{0.3ex}.


\section*{Résumé des commandes interactives \newline de pilotage du robot}

Les commandes de pilotage du \robot\ pour dessiner interactivement
avec le programme {\idt Des\-si\-ner} sont résumées dans la table \ref{tab:cmd-inter} 
(page \pageref{tab:cmd-inter}).

\begin{table}[htbp]
\begin{center}
\begin{tabular}{@{}|l|l|p{8.5cm}|@{}} \hline
{\em Touche\ :} & {\em Commande\ :} & {\em Effet\ :} \\ \hline  \hline
\multicolumn{3}{|l|}{\em Grille des points d'arrêt du \robot\ :} \\ \hline
\touche{.} & {\idt dg} & \underline{d}essine la \underline{g}rille des points d'arrêt du \robot. \\ \hline \hline
\multicolumn{3}{|l|}{\em Pivotements du \robot\ sur place :}\\ \hline
\touche{D} & {\idt phtd} & fait \underline{p}ivoter le \robot\ sur place d'un \underline{h}uitième de \underline{t}our à \underline{d}roite. \\
\touche{$+$} & {\idt pqtd} & fait \underline{p}ivoter le \robot\ sur place d'un \underline{q}uart de \underline{t}our à \underline{d}roite. \\
\touche{G} & {\idt phtg} & fait \underline{p}ivoter le \robot\ sur place d'un \underline{h}uitième de \underline{t}our à \underline{g}auche. \\
\touche{$-$} & {\idt pqtg} & fait \underline{p}ivoter le \robot\ sur place d'un \underline{q}uart de \underline{t}our à \underline{g}auche. \\ \hline \hline
\multicolumn{3}{|l|}{\em Déplacements du \robot\ :}\\ \hline
\touche{A} & {\idt av} & fait \underline{av}ancer le \robot\ tout droit jusqu'au prochain point d'arrêt du quadrillage. \\
\touche{)} & {\idt vd} & fait tourner le \robot\ par un \underline{v}irage d'un quart de cercle à \underline{d}roite. \\
\touche{(} & {\idt vg} & fait tourner le \robot\ par un \underline{v}irage d'un quart de cercle à \underline{g}auche. \\ \hline \hline
\multicolumn{3}{|l|}{\em Manœuvres du crayon du \robot\ :}\\ \hline
\touche{L} & {\idt lc} & \underline{l}ève le \underline{c}rayon du \robot, pour éviter de dessiner lors de ses déplacements. \\
\touche{B} & {\idt bc} & \underline{b}aisse le \underline{c}rayon du \robot, pour dessiner lors de ses déplacements. \\
\touche{P} & {\idt dp} & \underline{d}essine un \underline{p}oint, sans déplacer le \robot. \\
\touche{C} & {\idt dc} & \underline{d}essine une \underline{c}roix, sans déplacer le \robot. \\ \hline \hline
\multicolumn{3}{|l|}{\em Erreurs de pilotage du \robot\ :}\\ \hline
\multicolumn{2}{|l|}{%touche d'effacement 
\framebox[11.mm][l]{$\leftarrow$~~~}\hspace{0.3ex}
} & efface, dans l'ordre inverse, les derniers tracés du \robot. \\ \hline \hline
\multicolumn{3}{|l|}{\em Arrêt du \robot\ :}\\ \hline
\multicolumn{2}{|l|}{%retour a la ligne
\framebox[11.mm][l]{$\hookleftarrow$~~~}\hspace{0.3ex}
} & termine le dessin (qui peut ensuite être reproduit et imprimé). \\ \hline
\end{tabular}
\caption{Résumé des commandes interactives du \protect\robot\label{tab:cmd-inter}}
\end{center}
\end{table}

\clearpage\nopagebreak
\noindent \rule{137mm}{0.1mm} %---------------------------------------------
\begin{exercice}{
%\label{exo:trouver-9}
Supposons que le \robot\ soit initialement orienté vers le nord
avec son crayon baissé \robotCbN\/.  Parmi les chiffres d'une montre digitale 
re\-pré\-sen\-tés à la figure \ref{fig:alphabet} (page \pageref{fig:alphabet}), quel est
celui obtenu en donnant au robot la {\em séquence de commandes}%
\footnote{\ Une séquence \index{Séquence} de commandes est une succession ou suite de commandes données dans un certain ordre.}
  \hbox{ci-dessous~:}\\
\hspace*{2ex}{\idt pqtd$;$ av$;$ pqtg$;$ av$;$ av$;$ pqtg$;$ av$;$ pqtg$;$ av$;$ pqtg$;$ av$;$}\\
c'est-à-dire en tapant \hbox{successivement~:}\\
\hspace*{2ex}\touche{$+$} \touche{A} \touche{$-$} \touche{A} \touche{A} \touche{$-$} \touche{A} \touche{$-$} 
\touche{A} \touche{$-$} \touche{A}$\:$?
$\Box$}\end{exercice}

\begin{exercice}{
%\label{exo:trouver-F}
Quelle lettre majuscule de l'alphabet  de la figure \ref{fig:alphabet} (page \pageref{fig:alphabet})
obtient-on en donnant au \robot\ \robotCbN\ la séquence de commandes \hbox{suivante~:}\\
\hspace*{2ex}{\idt av$;$ av$;$ pqtd$;$ av$;$ pqtd$;$ pqtd$;$ av$;$ pqtd$;$ av$;$ av$;$ pqtd$;$ av$;$ av$;$}\\
c'est-à-dire en \hbox{tapant~:}\\
\hspace*{2ex}\touche{A} \touche{A} \touche{$+$} \touche{A} \touche{$+$} \touche{$+$} \touche{A} \touche{$+$} \touche{A} 
\touche{A} \touche{$+$} \touche{A} \touche{A}$\:$?
$\Box$}\end{exercice}

\begin{exercice}{
%\label{exo:faire-T}
Quelle est la suite de commandes à donner au \robot\ \robotCbN\ pour qu'il dessine 
la lettre {\tt T} majuscule représentée à la figure \ref{fig:alphabet} (page \pageref{fig:alphabet})$\:$?
$\Box$}\end{exercice}

\begin{exercice}{
%\label{exo:faire-4}
Quelle est la suite de commandes à donner au \robot\ \robotCbN\ pour qu'il dessine le chiffre 4 
représenté à la figure \ref{fig:alphabet} \hbox{(page \pageref{fig:alphabet})$\:$?}
$\Box$}\end{exercice}

\begin{exercice}{
%\label{exo:faire-octogone}
Quelle est la séquence de commandes à donner au \robot\ pour
dessiner l'octogone de la figure \ref{fig:octogone-cercles} 
\hbox{ci-dessous$\:$?}
$\Box$}\end{exercice}

\begin{exercice}{
%\label{exo:faire-cercles-tangents}
Quelle est la séquence de commandes à donner au \robot\ pour
dessiner les trois cercles tangents de la figure \ref{fig:octogone-cercles} 
\hbox{ci-dessous$\:$?}
$\Box$}\end{exercice}

\begin{figure}[htbp]
 \includehvspacePostScript[47mm,2.2cm]{octogone-cercles}{0.95}{0.95}
 \caption{Octogone et cercles tangents de tailles différentes}
 \label{fig:octogone-cercles}
\end{figure}
%%%%%%%%%%%%%%%%%%%%%%%%%%%%%%%%%%%%%%%%%%%%%%%%%%%%%%%%%%%%%%%%%%%%%%%%%%%%%%%%%%%%%%%%%%%

\begin{corrige}{
%\label{cor:trouver-9} 
En essayant sur l'ordinateur on trouve le chiffre neuf.
$\Box$}\end{corrige}

\begin{corrige}{
%\label{cor:trouver-F} 
La lettre F.
$\Box$}\end{corrige}

\begin{corrige}{
%\label{cor:faire-T} 
La lettre T majuscule de la figure \ref{fig:alphabet} (page \pageref{fig:alphabet}) est
obtenue en donnant successivement au \robot\ \robotCbN\ les commandes suivantes~:\\
\hspace*{2ex}{\idt av$;$ av$;$ av$;$ av$;$ pqtd$;$ av$;$ pqtd$;$ pqtd$;$ av$;$ av$;$}\\
c'est-à-dire en tapant \hbox{successivement~:}\\
\hspace*{2ex}\touche{A} \touche{A} \touche{A} \touche{A} \touche{$+$} \touche{A} \touche{$+$} \touche{$+$} 
\touche{A} \touche{A}.
$\Box$}\end{corrige}

\begin{corrige}{
%\label{cor:faire-4} 
Le chiffre 4 d'une montre digitale est
obtenu en donnant au \robot\ \robotCbN\ la séquence de commandes suivante~:\\
\hspace*{2ex}{\idt av$;$ av$;$ pqtd$;$ pqtd$;$ av$;$ pqtd$;$ av$;$ pqtd$;$ av$;$}\\
c'est-à-dire en tapant \hbox{successivement~:}\\
\hspace*{2ex}\touche{A} \touche{A} \touche{$+$} \touche{$+$} \touche{A} \touche{$+$} \touche{A} \touche{$+$} \touche{A}.
$\Box$}\end{corrige}

\begin{corrige}{
%\label{cor:faire-octogone} 
L'octogone de la figure \ref{fig:octogone-cercles} (page \pageref{fig:octogone-cercles}), 
est obtenu en donnant au \robot\ la séquence de commandes ci-dessous~:\\
\hspace*{2ex}{\idt av$;$ phtd$;$ av$;$ phtd$;$ av$;$ phtd$;$ av$;$ phtd$;$ av$;$ phtd$;$ av$;$}\\
\hspace*{2ex}{\idt phtd$;$ av$;$ phtd$;$ av$;$ }\\
c'est-à-dire en tapant \hbox{successivement~:}\\
\hspace*{2ex}\touche{A} \touche{D} \touche{A} \touche{D} \touche{A} \touche{D} 
\touche{A} \touche{D} \touche{A} \touche{D} \touche{A} \touche{D} 
\touche{A} \touche{D} \touche{A}.
$\Box$}\end{corrige}

\begin{corrige}{
%\label{cor:faire-cercles-tangents} 
Les cercles tangents de la figure \ref{fig:octogone-cercles} (page \pageref{fig:octogone-cercles}), 
sont obtenus en donnant au \robot\ la séquence de commandes \hbox{ci-dessous~:}\\
\hspace*{2ex}{\idt phtd$;$ }\\
\hspace*{2ex}{\idt vd$;$ vd$;$ vd$;$ vd$;$ vd$;$ vd$;$ }\\
\hspace*{2ex}{\idt phtg$;$ }\\
\hspace*{2ex}{\idt vg$;$ vg$;$ vg$;$ vg$;$ vg$;$ vg$;$ }\\
\hspace*{2ex}{\idt phtg$;$ }\\
\hspace*{2ex}{\idt vd$;$ vd$;$ vd$;$ vd$;$ }\\
c'est-à-dire en tapant  \hbox{successivement~:}\\
\hspace*{1ex}\touche{D} \\
\hspace*{1ex}\touche{)} \touche{)} \touche{)} \touche{)} \touche{)} \touche{)} \\
\hspace*{1ex}\touche{G} \\
\hspace*{1ex}\touche{(} \touche{(} \touche{(} \touche{(} \touche{(} \touche{(} \\
\hspace*{1ex}\touche{G} \\
\hspace*{1ex}\touche{)} \touche{)} \touche{)} \touche{)}. 
$\Box$}\end{corrige}

%%%%%%%%%%%%%%%%%%%%%%%%%%%%%%%%%%%%%%%%%%%%%%%%%%%%%%%%%%%%%%%%%%%%%%%%%%%
\noindent \framebox[\textwidth]{
\begin{minipage}{\myparboxwidth}
\setlength{\parindent}{2em}
\small \sf
\vspace{\abovedisplayskip}
\section*{Addenda}
Dans ce chapitre nous avons appris à piloter le \robot\ de manière
interactive.  Le {\em mode interactif} \index{Mode interactif} est le mode de fonctionnement de nombreux
programmes d'ordinateur que l'on appelle des {\em interpréteurs}. \index{Interpréteur}Ces programmes
engagent un dia\-lo\-gue avec l'utilisateur en utilisant un langage de
communication que l'on appelle le {\em langage de commande}\index{Langage de commande}. L'utilisateur
dirige le travail de l'interpréteur en lui donnant des commandes au moyen
du clavier de l'ordinateur, de la souris, etc. Les exemples sont très
nombreux. Citons les éditeurs de texte, les tableurs, le
`Finder' du Macintosh ou l'interpréteur de commandes des systèmes d'exploitation%
\footnote{\ Les systèmes d'exploitation sont des programmes qui gèrent l'organisation interne et l'affectation
des ressources des ordinateurs.
} (comme MS-DOS sur IBM-PC ou Unix sur les stations de travail scientifiques), les systèmes
transactionnels d'interrogation de bases de données (comme celles qui sont 
in\-ter\-ro\-gées par minitel). L'avantage évident
de cette interaction avec l'ordinateur est que l'utilisateur peut corriger
ses erreurs dès qu'elles se produisent. L'inconvénient majeur est que
les tâches répétitives deviennent très vite fastidieuses.

Pour notre \robot, le langage de commande se réduit à quelques
commandes très simples données au clavier (\hspace{-0.4ex}\touche{A}, \touche{B}, 
\touche{C}, \touche{D}, \touche{G}, \hbox{\touche{L},} \touche{P}, \touche{$+$}, 
\hbox{\touche{$-$},} \touche{(}, \touche{)}, \touche{.},  \touche{*}, 
\framebox[11.mm][l]{$\hookleftarrow$~~~}\hspace{0.3ex}) sans oublier la souris
ou, à défaut, les flèches de défilement qui sont des moyens très
rapides pour désigner un point précis sur l'écran.
Le \robot\ ne communique avec l'utilisateur que par des 
\raisebox{0.25ex}{$\scriptscriptstyle \boldmath\ll$}~%
bips-bips%
~\raisebox{0.25ex}{$\scriptscriptstyle \boldmath\gg$}
lui signalant 
qu'il ne doit pas le faire sortir du cadre. Les erreurs de dessin peuvent être 
corrigées par la touche d'effacement \framebox[11.mm][l]{$\leftarrow$~~~} 
qui annule l'effet des dernières commandes données au \robot.

Pendant le dessin interactif, le \robot\ engendre un programme \pascal\
qui peut ensuite être exécuté
pour reproduire et imprimer le dessin. Ceci peut être comparé aux robots
de peinture de l'industrie automobile. Un ouvrier spécialisé peint une
voiture au pistolet. Tous ses mouvements sont en\-re\-gis\-trés sur 
ordinateur puis utilisés pour engendrer un programme. En exécutant ce 
programme, l'ordinateur peut ensuite commander un robot disposant d'un bras 
manipulateur portant un pistolet pour reproduire fidèlement tous ces 
mouvements et peindre indéfiniment des voitures. Il s'agit d'un exemple rudimentaire
d'{\em apprentissage}, un thème de recherche actif en {\em intelligence
artificielle} où l'ordinateur engendre son propre 
programme pour réaliser certaines tâches à partir d'exemples qui lui sont fournis
en plus ou moins grand nombre.

\vspace{\abovedisplayskip}
\end{minipage}
}

\noindent \framebox[\textwidth]{
\begin{minipage}{\myparboxwidth}
\setlength{\parindent}{2em}
\small \sf
\vspace{\abovedisplayskip}
\indent Le dessin interactif sur ordinateur avec le \robot\ n'est pas plus difficile
et est certainement plus formateur que n'importe quel jeu vidéo. Il apprend
aux enfants à utiliser un {\em formalisme} simple. Il demande un effort
d'imagination pour s'orienter dans le plan puisque les commandes ne sont
pas relatives à l'observateur mais au \robot. 
Cependant la tâche
devient vite pénible pour les dessins géométriques répétitifs. 
Seule la programmation, que nous abordons dans le prochain
chapitre permet d'obtenir de tels dessins sans peine.

\vspace{\abovedisplayskip}
\end{minipage}
}
%%%%%%%%%%%%%%%%%%%%%%%%%%%%%%%%%%%%%%%%%%%%%%%%%%%%%%%%%%%%%%%%%%%%%%%%%%%%
